% ---
% id: protocol-compiler/rem:chaincond-reductions
% kind: rem
% title: Chain conditions here and in the reductions module
% summary: >-
%   The compiler's (C1)–(C3) are conditions on the record's fields and differ
%   from the reductions module's composability, tuple chaining and stage order,
%   whose numbers they share. That module proves the chain composes; this one
%   decides on the record what can be decided there and names the rest.
% prominence: standard
% status: accepted
% uses: [def:chaincond]
% mentions: [reductions/def:chain, reductions/def:link, reductions/thm:chain]
% assumes: []
% refs: []
% provenance: newly-derived
% lean:
% ---

\begin{remark}\label{rem:chaincond-reductions}
The two modules that meet at the chain record describe it from two sides: the
reductions module \emph{defines} the chain and \emph{proves} that it composes
(\xref{reductions/def:chain}, \xref{reductions/def:link},
\xref{reductions/thm:chain}); this module \emph{decides} what it can decide
on the presented record and \emph{names} what it cannot. Nothing is asserted
twice: the numbered conditions of Definition~\ref{def:chaincond} are this
module's, are stated on the record's fields alone, and are not the
composability (C1), tuple chaining (C2) and stage order (C3) of that module's
definition, whose numbering they happen to share --- ``(C1) of
Definition~\ref{def:chaincond}'' and ``reductions' composability (C1)'' are
different conditions on the same object.
\end{remark}
